% !TeX root = ../../Book1.tex
% 纲要标题：### A.3 积拓扑、紧性与投射有限空间
% 纲要位置：计划纲要.md，第 797 行。
% BEGIN APPENDIX OUTLINE SYNC
% #### A.3.1 积拓扑与连续映射
%
% #### A.3.2 有限交性质与紧性
%
% #### A.3.3 投射有限空间与有限层检测
% END APPENDIX OUTLINE SYNC

\section{积拓扑、紧性与投射有限空间}
\label{b1:sec:A03}

无限 Galois 群中的“接近”由有限层上取值相同来表达。这里只整理与这种有限信息有关的拓扑，不以度量、序列极限或完整点集拓扑作为前提。尤其对任意指标的积空间，基本开集只限制有限多个坐标。

\subsection{积拓扑与连续映射}
\label{b1:subsec:A03:01}

\begin{definition}\label{b1:def:topology-appendix}
集合 \(X\) 上的拓扑是一族包含 \(\varnothing,X\)、对任意并及有限交封闭的子集，族中成员称为开集，开集的补集称为闭集。包含一点的开集称为该点的开邻域。映射 \(f:X\rightarrow Y\) 连续，指 \(Y\) 中每个开集的原像在 \(X\) 中开。若两个不同的点总有不相交的开邻域，称空间为 Hausdorff 空间。
\end{definition}
子集 \(A\subseteq X\) 的子空间拓扑由所有 \(A\cap U\)（\(U\) 在 \(X\) 中开）组成；它的闭集恰为 \(A\cap F\)（\(F\) 在 \(X\) 中闭）。离散拓扑把每个子集都列为开集，因而离散空间中每个子集也闭。上述约定与\subsecref{b1:subsec:2202:02}一致。

\begin{definition}\label{b1:def:product-topology-appendix}
设 \((X_i)_{i\in I}\) 是拓扑空间族，\(P=\prod_iX_i\)，\(p_i:P\rightarrow X_i\) 是投影。积拓扑的基本开集为
\[
\bigcap_{i\in F}p_i^{-1}(U_i),
\qquad F\subseteq I\text{ 有限},\quad U_i\subseteq X_i\text{ 开}.
\]
所有这种集合的任意并构成积拓扑；空的 \(F\) 给出全空间。
\end{definition}
这些集合对有限交封闭，且覆盖 \(P\)，所以它们的任意并确实满足拓扑公理。若各 \(X_i\) 离散，可进一步以规定有限多个坐标的确切取值作为基本开集。特别地，\(\{0,1\}^{\mathbb N}\) 中规定所有坐标都为零得到一个闭单点，却不是开集：只固定有限坐标时总还能改动别的坐标。

\begin{proposition}\label{b1:prop:product-continuity}
映射 \(f:Y\rightarrow\prod_iX_i\) 连续，当且仅当每个坐标映射 \(p_i f\) 连续。若每个 \(X_i\) 为 Hausdorff 空间，则其积及积的任意子空间也为 Hausdorff 空间。
\end{proposition}
\begin{proof}
每个 \(p_i\) 的开集原像是基本开集，所以连续，得到必要性。反之，基本开集在 \(f\) 下的原像是有限多个 \((p_i f)^{-1}(U_i)\) 的交，故开；任意开集是基本开集的并，其原像仍开。两个不同的积空间点在某个坐标上不同；在该坐标取不相交开邻域，其投影原像将两点分离。子空间中的分离邻域只需与子空间相交。
\end{proof}
因此，\cref{b1:prop:inverse-limit-set-universal}中的极限若采用积的子空间拓扑，且各 \(f_i\) 连续，则诱导映射也连续。这说明逆极限的拓扑无需另行猜测：它就是要求全部坐标连续所产生的拓扑。

\subsection{有限交性质与紧性}
\label{b1:subsec:A03:02}

沿用正文约定：每个开覆盖有有限子覆盖称为拟紧；拟紧且 Hausdorff 才称为紧。这一点须与把“紧”直接用作“拟紧”的文献区别。集合族 \(\mathcal F\) 具有有限交性质，指任意有限个成员的交非空，包括零个成员；空交取为全空间。

\begin{proposition}\label{b1:prop:compact-fip}
空间 \(X\) 拟紧，当且仅当每个具有有限交性质的闭集族总交非空。拟紧空间的闭子空间仍拟紧；紧空间的闭子空间仍紧。
\end{proposition}
\begin{proof}
闭集族 \((F_\alpha)\) 总交为空，等价于其开补集覆盖 \(X\)；其中某个有限交为空，等价于相应有限多个开补集已经覆盖 \(X\)。将拟紧性的有限子覆盖条件取逆否命题，正好得到所述等价性。

若 \(C\subseteq X\) 闭，则 \(C\) 中的闭集也是 \(X\) 中的闭集。其有限交非空可以在 \(X\) 中检验，由刚证的判据得到总交非空，所以 \(C\) 拟紧。Hausdorff 性对子空间成立，故第二个结论也成立。空子空间的开覆盖直接有空的有限子覆盖，同样符合结论。
\end{proof}

\begin{proposition}\label{b1:prop:compact-maps}
拟紧空间在连续映射下的像仍拟紧。Hausdorff 空间中的拟紧子空间是闭集。因此，从紧空间到 Hausdorff 空间的连续双射是同胚。
\end{proposition}
\begin{proof}
像空间的开覆盖拉回给出源空间的开覆盖，有限子覆盖再映回覆盖整个像，得到第一句。设 \(C\subseteq Y\) 拟紧，\(Y\) 为 Hausdorff 空间，\(y\notin C\)。对每个 \(c\in C\)，选不相交的开邻域 \(U_c\ni c\)、\(V_c\ni y\)。\(U_c\) 覆盖 \(C\)，取有限子覆盖 \(U_{c_1},\ldots,U_{c_r}\)。开邻域 \(V_{c_1}\cap\cdots\cap V_{c_r}\) 与 \(C\) 不交，所以 \(Y\setminus C\) 开，\(C\) 闭；\(C=\varnothing\) 时直接成立。

若 \(f:X\rightarrow Y\) 为题述连续双射，\(X\) 的闭集由\cref{b1:prop:compact-fip}拟紧，其像也拟紧，进而在 \(Y\) 中闭。这表示 \(f^{-1}\) 的每个闭集原像都闭，等价于 \(f^{-1}\) 连续，故为同胚。
\end{proof}
本卷所需的积紧性是\cref{b1:lem:finite-discrete-product-compact}：任意有限离散空间族的积是紧空间。其完整证明在\subsecref{b1:subsec:2202:04}，从有有限交性质的闭集族生成真滤子，用 Zorn 引理扩张为极大真滤子，再在每个有限坐标上取得唯一被选中的纤维。有限交判据给出总交中的点。这里汇总该结论，不把正文中的证明改成对附录的省略引用，也不在此额外假设一般的积紧性定理。

\subsection{投射有限空间与有限层检测}
\label{b1:subsec:A03:03}

\begin{definition}\label{b1:def:profinite-space}
与某个有限离散集合的有向逆极限同胚的空间，称为\term[toushe-youxian-kongjian]{投射有限空间}[profinite space]。它的拓扑是相容族在各有限层直积中的子空间拓扑。若同时具有逐坐标的群结构且过渡映射都是群同态，就得到正文的投射有限群。
\end{definition}

\begin{proposition}\label{b1:prop:finite-inverse-limit}
设 \((X_i,r_{ji})\) 是非空有向集上的有限离散集合逆系统，\(X=\varprojlim_iX_i\)。则 \(X\) 紧，且其基本开集可取为 \(\pi_i^{-1}(\{a\})\)，这些集合同时闭。因而 \(X\) 全不连通。若每个 \(X_i\) 非空，则 \(X\) 非空；若过渡映射还全都满射，则每个投影 \(\pi_i:X\rightarrow X_i\) 满射。
\end{proposition}
\begin{proof}
在 \(P=\prod_iX_i\) 中，条件 \(r_{ji}(x_j)=x_i\) 定义闭集 \(C_{ji}\)：不满足条件的每个点都由这两个坐标的取值给出完全落在补集中的基本开邻域。因此 \(X=\bigcap_{i\leq j}C_{ji}\) 闭。\cref{b1:lem:finite-discrete-product-compact}和\cref{b1:prop:compact-fip}给出 \(X\) 紧。

\(X\) 中一个基本开邻域原本规定有限个坐标 \(i_1,\ldots,i_r\)。取共同上层 \(k\)，固定点的第 \(k\) 坐标便同时固定这些坐标，所以单层纤维构成邻域基。纤维的补集是同层其他纤维之并，故它也闭。两个不同点在某坐标不同，相应开闭纤维将任何同时包含它们的子空间分离为两个非空相对开部分；因此非空连通子空间只能是单点，即全不连通。

现设各层非空。采用选择公理先在 \(P\) 中取一点作为未规定坐标的默认值。给定有限多个相容条件，取高于其中全部指标的 \(k\)，任选 \(a_k\in X_k\)。将涉及的坐标 \(i\) 赋值为 \(r_{ki}(a_k)\)，其余坐标用默认值填充。复合条件保证这有限多个相容条件同时成立。所以闭集族 \((C_{ji})\) 具有有限交性质；由积的紧性和\cref{b1:prop:compact-fip}，总交 \(X\) 非空。

最后固定 \(i\) 和 \(a\in X_i\)，在上述闭集族中加入条件 \(x_i=a\)。处理任意有限个条件时，把共同上层 \(k\) 也取在 \(i\) 之上。过渡映射满射使我们能选 \(a_k\) 满足 \(r_{ki}(a_k)=a\)，于是同一有限交论证产生 \(X\) 中第 \(i\) 坐标为 \(a\) 的点，证明投影满射。
\end{proof}
非空结论与投影满射结论的假设不同。即使每层都是 \(\{0,1\}\)，若严格后层到前层的映射恒为零，极限也只有全零族，投影并不满射。有限群的逆极限总有单位族，故其非空性不需要这个存在性论证；满射问题仍要单独处理。

\begin{proposition}\label{b1:prop:finite-value-factorization}
设 \(X=\varprojlim_iX_i\) 如上，\(D\) 为有限离散集合，\(f:X\rightarrow D\) 连续。则存在指标 \(k\) 及映射 \(\bar f:\pi_k(X)\rightarrow D\)，使 \(f=\bar f\pi_k\)。若 \(\pi_k\) 满射，\(\bar f\) 定义在整个 \(X_k\) 上。
\end{proposition}
\begin{proof}
若 \(X\) 为空，取任一指标并用空映射即可。否则，每个 \(x\in X\) 都有一个单层纤维邻域 \(U_x\)，使 \(f\) 在该邻域上恒等于 \(f(x)\)，因为 \(f^{-1}(\{f(x)\})\) 开。由\cref{b1:prop:finite-inverse-limit}的紧性，有限多个 \(U_{x_1},\ldots,U_{x_r}\) 覆盖 \(X\)。取高于它们所用指标的共同上层 \(k\)。

若 \(\pi_k(y)=\pi_k(z)\)，选一个含 \(y\) 的 \(U_{x_j}\)。相容性说明 \(y,z\) 在定义此邻域的坐标上相同，所以 \(z\) 也属于它，进而 \(f(y)=f(z)\)。由\cref{b1:prop:set-quotient-universal}，\(f\) 唯一下降到 \(\pi_k(X)\)，得到所需分解。
\end{proof}
例如令 \(X_n=\mathbb Z/2^n\mathbb Z\)、\(n\geq1\)，过渡映射为约化。相容族恰好由一列数字 \(\varepsilon_0,\varepsilon_1,\ldots\in\{0,1\}\) 描述：第 \(n\) 个坐标是
\[
\varepsilon_0+2\varepsilon_1+\cdots+2^{n-1}\varepsilon_{n-1}\pmod{2^n}.
\]
给定相容族，先从模二坐标读出 \(\varepsilon_0\)；已读出前 \(n\) 位后，相容性使模 \(2^{n+1}\) 的代表元与此前部分和之差唯一为 \(0\) 或 \(2^n\)，从而确定下一位。这证明两种数据互相唯一决定。这里的“无限和”仅指全部有限截断组成的相容族，没有使用实数中的级数收敛。连续的有限值观察只能依赖有限多个起始数字；整个空间却含有无限多个点，而且没有孤立点。
